\section{Introduction}
\label{sec:introduction}

Self-improving agents work today wherever the supervision signal is dense. Give a
system a kernel to make faster, a leaderboard to climb, or a unit test to turn green,
and iteration has something to optimize: every attempt returns a number, and the
number says whether to keep going. Progress under those conditions is real, and it is
also narrow. It describes a student who has become excellent at answering questions
someone else has already written down.

The tasks that motivate this report supply no such signal. A mathematical campaign
rarely terminates in the theorem it set out to prove: intermediate bounds,
counterexamples, and reformulations are the normal output, and they carry real value.
A software request is often underspecified until a candidate implementation exposes
what was missing. In system verification, the specification and the implementation may
both be wrong, and deciding which one to change is itself part of the work. In chip
design and in materials research, the measurement that would settle the question is
frequently the thing under construction. What these share is not difficulty but
\emph{underdefinition}: at the moment work starts, nobody, human or machine, can state
the objective precisely enough to optimize against it.

This changes what a runtime is for. If the score cannot be trusted to steer, the
system's job is not to climb it but to hold a campaign together long enough for the
objective itself to become clear, and to route the specific contradictions it cannot
resolve to the human expert who can. We call this \emph{expert--agent co-evolution}:
the standing intent stays stable while the operational objective, constraints, and
verification criteria are refined by evidence produced along the way, and the moments
that require human authority are surfaced explicitly rather than guessed past.

Permitting that revision is difficult because it is externally indistinguishable from
failure. A system that abandons its stated target may have discovered that the target
was misspecified, or it may have failed and then rationalized. Nothing in the final
artifact separates the two. A runtime that allows objective revision without resolving
this ambiguity invites \emph{goal drift}: the objective degrades toward whatever the
executor happens to be able to finish. The standard response is to remove the choice,
fixing the objective and optimizing execution toward it.

Prior work follows that response. ReAct, SWE-agent, and OpenHands established
practical interfaces between models and external environments
~\citep{yao2023react,yang2024sweagent,wang2024openhands}. The AI Scientist,
CycleResearcher, AutoSci, FARS, and Arbor extend this interaction into iterative
research programs covering planning, experiments, review, memory, and scientific
communication~\citep{lu2024aiscientist,weng2025cycleresearcher,qian2026autosci,
tang2026fars,jin2026arbor}. These systems differ in how they search, remember, and
review, but they share the assumption that the objective is given. Their reported
progress makes the remaining problem visible: agent success declines as software-task
horizons grow~\citep{kwa2025longtasks}, and the horizons that matter most are those
over which the objective is least likely to survive unchanged.

We take the opposite position: pivoting should be permitted and made safe by
verification. An admissible pivot is supported by evidence that the previous route
or objective was unreachable or misspecified, admitted through an explicit role
boundary, and recorded so that later missions inherit both the change and its
justification. Verification is therefore not a quality filter applied after
execution. It is the mechanism that makes pivoting distinguishable from drift.

Three properties must hold before any such certification is possible, and each fails
in long-running agents. \emph{Continuity} fails when a growing transcript is compacted
or dropped, so the evidence that motivated a revision no longer exists when the
revision is proposed. \emph{Acceptance} fails when the component that performs the
work also declares it complete, so a revision is certified by the party with the
strongest incentive to accept it. \emph{Experience} fails when only the final artifact
survives, so refuted routes cannot later be cited as evidence that an objective is
unreachable. Hierarchical and workflow-oriented memory systems address parts of this
problem~\citep{packer2023memgpt,wang2024workflowmemory,xu2025amem}. A larger context
window extends session continuity; certifying a change of objective additionally
requires explicit state, ownership, and update rules.

We present \systemname (Autonomous Research Generation and Understanding System), a
general-purpose agentic runtime for long-horizon reasoning. Its central abstraction
is a sequence of \emph{bounded missions} executed against a durable project state.
Four model-driven roles divide authority: a Manager anchors the objective and
campaign state, a Planner selects the next units of work, an Engineer implements
and evaluates them, and a Reviewer inspects artifacts and run outputs before issuing
the completion verdict when independent review is required or requested. Allowed
low-risk tasks may instead use recorded Engineer self-review. These roles operate across three planes---control,
execution, and records---so that scheduling, work, and record keeping remain
separable even though they participate in one continuous loop.

At the report level, we summarize the evolving working contract as
$K_t=(\iota,o_t,c_t,v_t)$. Here $\iota$ is the standing user intent that the campaign
must preserve; $o_t$ is the current operational objective; $c_t$ is the set of known
constraints; and $v_t$ specifies the verification criteria at mission $t$. We use
$X_t$ for user-visible clarifications, priorities, and unresolved questions. This
notation separates an evidence-backed refinement of the operational contract from a
silent change of intent. Section~\ref{sec:problem} formalizes how evidence, recorded
verdicts, and the required user or Manager authority admit a material update.

Deployment provides qualitative motivation for this design. Expert users of \systemname
report that a substantial part of its practical value lies in where it stops: the
system declines to continue on an underspecified objective, and that refusal surfaces
constraints the user had not stated. Related internal system-verification cases
involve real researchers and project-specific details that cannot be disclosed in
this report; we therefore record the observed use as qualitative motivation rather
than a public prospective user-study result. After initial assignment, the current runtime can execute
ordinary rounds unattended and pause explicitly for operator-owned decisions; the
public traces do not provide a measured zero-touch rate.

A pivot should not reset the campaign. The runtime self-evolves the approach between
the standing intent and the current operational contract: admitted Stages retain
attempted routes, measurements, verdicts, accepted and rejected results, reusable
skills, verifiers, and routing decisions. Candidate updates enter persistent state
only after the authorized role checks the relevant artifacts and task-native evidence;
the gate may include an official verifier, independent Reviewer, or explicitly allowed
Engineer self-review according to task policy. Model parameters remain fixed, but
later missions therefore begin from a changed search policy rather than merely a
longer transcript. We call this mechanism \emph{verification-gated fixed-model runtime
self-evolution}. Section~\ref{sec:results} reports its longitudinal evidence as
observational rather than a controlled learning ablation.

We sharpen this design language into a process-to-capability theory. The final
artifact is a projection of a richer typed trajectory; reviewed compression retains
the parts needed for later decisions; and an accepted update counts as compounding
only when it reduces future task risk or resource cost. This separates three claims
that are often conflated: more process information is available, a review gate
improves admission quality, and retained state actually helps later work.

Overall, the evaluation asks whether one self-evolving runtime can combine broad
task capability, verified accumulation across missions, and sustained research
delivery. The evidence is organized in three parts. First, seven benchmark arenas
cover software repair, GPU kernels, language-model training, training speed,
research-assistant tasks, and mathematical data synthesis. Beyond SWE-Bench Pro,
the suite includes 76.8\% on AARRI-Bench and a 28.0 mathematical-data gap, both above
their reported references; a de-duplicated inventory separately records 41 research
artifacts across six programs. Second, the 731-task SWE-Bench Pro trajectory supports
longitudinal analyses of Skill/Wiki evolution, verifier recovery, Reviewer
intervention, and mature-stage operating cost. Third, vertical evidence tests whether
retained state changes later research: one mathematical campaign preserves a
falsified route and six proof-backed theorem-frontier updates, while six
paper-production campaigns span \PaperCaseMissions{} missions and
\PaperCaseRollbacks{} Stage rollbacks before producing final AAAI- and ACL-formatted
manuscripts. Two further verticals carry the same discipline into domains with
unforgiving external checkers: a chip certified for its demonstrated scope through
mapped synthesis and static timing, and a materials campaign that replaces a
published sampling method with a simpler one after removing a confound in the
original comparison.

This report makes seven contributions:

\begin{enumerate}[label=\textbf{C\arabic*.}]
  \item We formulate long-horizon research as verified pivoting over a compact
    report-level contract $K_t=(\iota,o_t,c_t,v_t)$ and explicit user decision state
    $X_t$, with each variable defined above. The analytical \code{ManagerAdmit}
    operator makes evidence, authority, and
    provenance requirements explicit without claiming one atomic implementation API.
  \item We present a runtime in which contract refinement is represented through
    typed, logged, role-owned surfaces rather than one atomic operation or an implicit
    consequence of replanning. \code{GoalContract} revisions, Manager stage
    transitions, abandoned approaches, and recorded verdict sources give material changes
    an author, a precondition, and an audit trail. The same admission discipline
    gates fixed-model runtime self-evolution: generated memories, skills, verifiers,
    routing decisions, and rejected routes become reusable only after evidence checks
    and an authorized commit.
  \item We measure where the runtime spends verification and where it refuses to
    terminate. Of \ReviewerTasks{} SWE-Bench Pro tasks, \ReviewerInvoked{} invoke an
    independent Reviewer and \ReviewerSkipped{} use Engineer self-review. The Reviewer
    withholds completion on \ReviewerRevisionRequested{} tasks, of which
    \ReviewerVerifierRecovered{} later pass the official verifier and
    \ReviewerStrictRescues{} complete the strict review loop; a further
    \ReviewerFirstBlocked{} tasks are declared blocked rather than reported complete.
    A longitudinal view of the same run shows mature Waves using
    \SWEProTokenReduction\% fewer solve input Tokens and \SWEProTimeReduction\% less
    active workflow time per task than startup.
  \item We present a six-project autonomous paper-production case study covering
    \PaperCaseCampaignHours{} campaign-hours, \PaperCaseMissions{} bounded missions,
    \PaperCaseRounds{} Engineer rounds, \PaperCaseContinues{} Reviewer revisions,
    \PaperCaseSessionRolls{} session rolls, and \PaperCaseRollbacks{} Stage rollbacks,
    with all six canonical pipelines reaching submission completion. A representative
    campaign drops seven weak approaches before reframing a failed method search as a
    4,500-row negative-results study, then repairs two late submission defects
    without resetting the research state.
  \item We reconstruct one mathematical campaign as a role-resolved Agent trace,
    making the Manager, Planner, Engineer, and Reviewer actions visible across
    multiple bounded missions, and retaining one falsified route alongside six
    accepted theorem-frontier updates.
  \item We report a capability floor across seven benchmark arenas in one unified
    table that preserves each benchmark's backbone, backend, protocol, native metric,
    and reference, establishing that the verification machinery does not cost
    end-task performance: on SWE-Bench Pro, \systemname reaches approximately
    \SWEProArgusAccuracy\% accuracy versus \SWEProDirectAccuracy\% for Direct Copilot
    at \SWEProTokenRatio$\times$ aggregate Tokens.
  \item We report external adoption of an \systemname-produced artifact: a TileLang
    \code{RWKV6} kernel reviewed by an outside maintainer and merged into
    \code{fla-org:main}.
\end{enumerate}

The remainder of the report follows a conventional research-paper organization.
Section~\ref{sec:related} positions the system against prior agent and autonomous
research work. Section~\ref{sec:problem} formalizes the operating problem;
Section~\ref{sec:method} presents the runtime; Sections~\ref{sec:methodology} and
\ref{sec:results} describe the evaluation protocol and results; and
Sections~\ref{sec:discussion}--\ref{sec:conclusion} discuss interpretation,
limitations, and conclusions. The appendix contains detailed experimental tables
and notation.
